% Standalone theorem insert for W33 papers.
% File: paper/w33_flag_overlap_association_insert.tex

\section{The W33 Flag Overlap Association Theorem}
\label{sec:w33-flag-overlap-association}

The minimal support theorem identifies the projective minimal $X$ supports with
isotropic line-stars.  Equivalently,
\[
        X_{\min}\;\cong\;\{(p,L): p\in L,\; L\text{ a totally isotropic line of }W(3,3)\}.
\]
Thus $X_{\min}$ is the flag set of the generalized quadrangle $W(3,3)$.
Since $W(3,3)$ has $40$ points, $40$ lines, and $4$ points per line, this gives
\[
        |X_{\min}|=40\cdot4=160.
\]

\begin{theorem}[W33 Flag Overlap Association Theorem]
Let $f=(p,L)$ and $g=(q,M)$ be two distinct point-line flags of $W(3,3)$.
Let $\Omega(f,g)$ denote the number of minimal $Z$-quadrangles visible from both
minimal $X$ line-stars $f$ and $g$.  Then $\Omega(f,g)$ is determined exactly by
the relative position of the two flags:
\[
\begin{array}{c|c}
\text{relative position of }(p,L),(q,M) & \Omega(f,g)\\
\hline
p=q\text{ or }L=M & 27\\
\text{exactly one cross-incidence }(p\in M\text{ or }q\in L) & 9\\
L\cap M\ne\varnothing\text{ away from }p,q\text{, or }L,M\text{ skew and }p\sim q & 3\\
L,M\text{ skew and }p\not\sim q & 1.
\end{array}
\]
Consequently, for every fixed flag $f$, the overlap distribution over the other
$159$ flags is
\[
        27^6,\qquad 9^{18},\qquad 3^{54},\qquad 1^{81}.
\]
Equivalently,
\[
        1^{81},\qquad 3^{54},\qquad 9^{18},\qquad 27^6.
\]
The global unordered distribution is
\[
        1:6480,
        \qquad
        3:4320,
        \qquad
        9:1440,
        \qquad
        27:480.
\]
\end{theorem}

\begin{proof}[Finite certificate]
The certificate reconstructs $W(3,3)$ from the symplectic form on the projective
points of $\mathbb F_3^4$, enumerates the $160$ flags, classifies every ordered pair
of distinct flags by relative incidence position, and checks that each relative
position has a single overlap value.  The ordered relation counts are
\[
\begin{array}{c|c|c}
\text{relative position} & \text{ordered count} & \text{overlap}\\
\hline
p=q & 480 & 27\\
L=M & 480 & 27\\
p\in M\text{ only} & 1440 & 9\\
q\in L\text{ only} & 1440 & 9\\
L\cap M\ne\varnothing\text{ away from }p,q & 4320 & 3\\
L,M\text{ skew and }p\sim q & 4320 & 3\\
L,M\text{ skew and }p\not\sim q & 12960 & 1.
\end{array}
\]
These ordered counts sum to $160\cdot159$.  Dividing by $2$ gives the unordered
counts stated above.  Therefore the $3$-adic overlap value is a flag-distance
invariant on $W(3,3)$.
\end{proof}

\begin{corollary}[Geometric meaning of the $3$-adic scheme]
The unsigned minimal logical Gram matrix $UU^T$ is the adjacency algebra of a
point-line flag association scheme.  The four off-diagonal values
\[
        1,3,9,27
\]
are not arbitrary eigenvalue artifacts: they measure the relative incidence distance
between two local line-star defects in $W(3,3)$.
\end{corollary}
